\begin{textsample}{POS Dim 2 – sp – Score 29.00 – sp/sp\_em\_3.txt}  \label{ex:f2_pos_201}
Uma breve retrospectiva das discussões curriculares no Brasil e no Estado de São Paulo Segundo a história da educação brasileira, o Ensino Médio surge com as escolas secundárias, quando o projeto dos jesuítas dos séculos XVI e XVII implementa sistemas educacionais de dois níveis : o primário e o secundário.
Nesse mesmo período, o governo central criou um modelo de educação cuja meta era propiciar o ingresso dos estudantes brasileiros nas universidades europeias e, assim, garantir a formação de um grupo privilegiado da população.
A expansão do Ensino Médio, de acordo com Tartuce et al. ( 2015 ), orientou -se pela separação entre o Ensino Médio Regular e o Profissionalizante, sendo este construído sob o escopo da desvalorização social. “ Essa dualidade histórico-estrutural, acentuada pelas desigualdades socioeconômicas do país ” ( TARTUCE et al., 2015, p. 19 ), determinou a segmentação das \textbf{trajetórias} dos jovens estudantes entre a formação acadêmica – \textbf{destinada} às elites – e a formação profissional/técnica, reservada à classe \textbf{trabalhadora}.
No início do século XX, a \textbf{oferta} do Ensino Médio circunscreveu -se a estabelecimentos como os Liceus, nas capitais dos Estados, direcionados para a educação \textbf{masculina}, e às escolas normais, para o gênero feminino.
Segundo Zotti et al. ( 2015 ), essas escolas eram restritas às elites burocráticas e latifundiárias do país.
Somente a partir da década de 1930 principia -se uma \textbf{política} governamental oficial norteada à estruturação do sistema educacional brasileiro.
Na década de 1940, com a promulgação da Lei Orgânica do Ensino Secundário de 1942 ( BRASIL, 1942 ), ocorre a divisão entre formação geral e formação específica, ao estabelecer -se a divisão em dois ciclos, conforme os artigos 2º, 3º e 4º da lei : > Art. 2º O ensino secundário será ministrado em dois ciclos.
O primeiro compreenderá um só \textbf{curso} : o \textbf{curso} ginasial.
O segundo compreenderá dois \textbf{cursos} paralelos : o \textbf{curso} clássico e o \textbf{curso} científico. > > Art. 3º O \textbf{curso} ginasial, que terá a duração de quatro anos, \textbf{destinar} -se-á a dar aos \textbf{adolescentes} os elementos fundamentais do ensino secundário. > > Art. 4º O \textbf{curso} clássico e o \textbf{curso} científico, cada qual com a duração de três anos, terão por objetivo consolidar a educação ministrada no \textbf{curso} ginasial e bem assim desenvolvê -la e aprofundá -la.
No \textbf{curso} clássico, \textbf{concorrerá} para a formação intelectual, além de um maior conhecimento de filosofia, um acentuado estudo das letras antigas ; no \textbf{curso} científico, essa formação será marcada por um estudo maior de ciências.
Após esse período, foi sancionada a Lei de \textbf{Diretrizes} e Bases da Educação ( BRASIL, Lei nº 4.024/61 ), que equipara o Ensino Secundário Regular ao Ensino Profissionalizante, para fins de acesso ao Ensino Superior.
Assim, os egressos do 2º ciclo, de qualquer modalidade, poderiam seguir seus estudos acadêmicos.
Porém, nesse período era obrigatório somente o Ensino Primário, com quatro anos de duração, o que restringia o acesso ao Nível Secundário, havendo a exigência do exame de admissão para o acesso ao Ginásio, tornando mínima a parcela da população que \textbf{frequentava} até mesmo a escola primária.
Portanto, o acesso às escolas de Ensino Secundário era muito restrito.
No início dos anos 1970, foi promulgada a Lei de \textbf{Diretrizes} e Bases de 1971 – LDB 5.692/71 ( BRASIL, 1971 ), que definiu o direito social à Educação Básica, com a fusão do \textbf{curso} primário com o ciclo ginasial, tornando obrigatória a \textbf{escolaridade} de oito anos.
Nesse mesmo período é extinto o exame de admissão, permitindo que o estudante avance nos estudos.
Além disso, a lei traz a organização curricular com um núcleo comum e uma parte diversificada : > Art. 4º Os \textbf{currículos} do ensino de 1º e 2º graus terão um núcleo comum, obrigatório em âmbito nacional, e uma parte diversificada para \textbf{atender}, conforme as necessidades e possibilidades concretas, às peculiaridades locais, aos planos dos estabelecimentos e às diferenças individuais dos alunos.
Após uma década da promulgação da Lei nº 5.692/71, o Brasil reconhece as fragilidades dessa diretriz.
Com a promulgação da Lei nº 7.044/82 ( BRASIL, 1982 ), a \textbf{habilitação} \textbf{profissional} passa a ser uma opção do estabelecimento de ensino ; portanto, retira -se a obrigatoriedade da formação \textbf{profissional} ( ZIBAS et al., 2015 ).
A Constituição Federal de 1988 traz um marco importante para a educação nacional.
Foi a partir dela que a educação tornou -se um direito universal, conforme explicitado no artigo 205 : > A educação, direito de todos e dever do Estado e da família, será promovida e incentivada com a colaboração da sociedade, visando ao pleno desenvolvimento da pessoa, seu preparo para o exercício da cidadania e sua \textbf{qualificação} para o trabalho. ( BRASIL, 1988 ) Na década de 1990, as disposições legais sobre o Ensino Médio reconhecem a importância de uma sólida educação geral para levar o estudante a se constituir cidadão, fornecendo -lhes os meios para progredir no trabalho e/ou nos estudos.
Posteriormente, alicerçada na concepção de uma educação comum necessária a todos os cidadãos, a Lei de \textbf{Diretrizes} e Bases da Educação Nacional de 1996 – LDB 9.394/96 ( BRASIL, 1996 ) inclui o Ensino Médio como parte integrante da Educação Básica, isto é, como a formação essencial – científica, tecnológica, cultural e social – a ser garantida a todos os jovens, promovendo a possibilidade da continuidade dos estudos e/ou do ingresso no mercado de trabalho, por meio de uma sólida preparação para o mundo \textbf{profissional}.
Nos anos 2000, são publicados os Parâmetros Curriculares Nacionais para o Ensino Médio – PCNEM ( BRASIL, 2000 ), pontuando a dimensão de preparação para o trabalho e competências cognitivas, psicomotoras e socioafetivas que capacitam o estudante para a vida.
Elaborados em consonância com as indicações da LDB/96 e das \textbf{Diretrizes} Curriculares Nacionais para o Ensino Médio – DCNEM/98, os PCNEM apresentam as bases legais que fundamentam o então chamado “ Novo Ensino Médio ”, especificando as competências e habilidades por área de conhecimento e pelas disciplinas potenciais de cada área.
No ano de 2012, o Ministério da Educação ( MEC ) define novas \textbf{Diretrizes} Curriculares Nacionais para o Ensino Médio ( DCNEM/2012 ), apresentando os princípios e procedimentos a serem seguidos pelos sistemas de ensino e escolas, promovendo grandes mudanças diante de novas exigências educacionais resultantes das transformações do século XXI, e considerando a diversidade das juventudes e seus interesses.
Essas \textbf{diretrizes} corroboram as \textbf{legislações} anteriores, no que diz respeito à apresentação de um \textbf{currículo} composto por uma base nacional comum, à qual todos os estudantes devem ter acesso, e uma parte diversificada, a ser definida em cada sistema de ensino, constituindo assim um todo integrado.
Em 2017, a Lei de \textbf{Diretrizes} e Bases da Educação Nacional sofreu alterações pela Lei 13.415/2017.
Em seu \textbf{Capítulo} II, que trata da Educação Básica, foram alterados os artigos 24, 26, 35 -A e 36.
Os artigos 35 -A e 36 compõem a Seção IV do \textbf{Capítulo} II, que trata especificamente do Ensino Médio e prevê mudança na estrutura curricular, conforme apresentado no artigo 36 : > O \textbf{currículo} do ensino médio será composto pela Base Nacional Comum Curricular e por \textbf{itinerários} \textbf{formativos}, que deverão ser organizados por meio da \textbf{oferta} de diferentes arranjos curriculares, conforme a relevância para o contexto local e a possibilidade dos sistemas de ensino, a saber : > > I - linguagens e suas tecnologias ; > II - matemática e suas tecnologias ; > III - ciências da natureza e suas tecnologias ; > IV - ciências humanas e sociais aplicadas ; > V - formação \textbf{técnica} e \textbf{profissional}.
Essas alterações trouxeram a necessidade de atualização das DCNEM, que foram modificadas pela Resolução nº 3 de 21/11/2018. \textbf{Alinhadas} com as mudanças previstas na LDB, em seu art. 10, as DCNEM definem que a formação geral básica e o \textbf{itinerário} \textbf{formativo} do Ensino Médio devem ser partes indissociáveis.
Além disso, o artigo 12 apresenta como deve ser a organização dos \textbf{itinerários} \textbf{formativos}, a partir das áreas do conhecimento e da formação \textbf{técnica} e \textbf{profissional}.
No inciso 1º, apresenta -se : > Os \textbf{itinerários} \textbf{formativos} devem considerar as demandas e necessidades do mundo contemporâneo, estar \textbf{sintonizados} com os diferentes interesses dos estudantes e sua inserção na sociedade, o contexto local e as possibilidades de \textbf{oferta} dos sistemas e instituições de ensino. ( BRASIL, 2018 ) Essas mudanças foram estabelecidas visando responder às constantes transformações sociais e garantir a permanência e \textbf{qualidade} do processo de ensino e aprendizagem das instituições escolares.
Com isso, propõem a elaboração de \textbf{políticas} públicas voltadas a \textbf{atender} as necessidades dos jovens estudantes, em suas diferentes especificidades, próprias desta etapa de Educação Básica.
Outros marcos legais que fundamentam a construção da etapa do Ensino Médio do \textbf{Currículo} Paulista são a Base Nacional Comum Curricular ( BNCC ), \textbf{homologada} em 14 de dezembro de 2018, para a etapa do Ensino Médio, e os Referenciais Curriculares para a Elaboração de \textbf{Itinerários} \textbf{Formativos}, por meio da \textbf{Portaria} nº 1.432 de 28/12/2018.
Com base nesses dois marcos legais e em escutas à \textbf{rede} de ensino, foram construídos os organizadores curriculares das quatro áreas do conhecimento na formação geral básica e nos \textbf{itinerários} \textbf{formativos} e da formação \textbf{técnica} e \textbf{profissional} para o \textbf{itinerário} \textbf{formativo}.
Para tanto, faz -se necessário buscar caminhos para novas formas de organização escolar visando resolver os problemas da defasagem idade-série e da evasão dos estudantes do Ensino Médio.
Acreditamos que só por meio de ações efetivas é possível promover o tripé “ acesso, permanência e \textbf{qualidade} social ” da educação, como também as relações dos jovens e adultos com o mundo do trabalho e a empregabilidade.
É sob as premissas da autonomia e da flexibilização que a Lei nº 13.415/2017 traz, em sua essência, possibilidades de fomentar ao jovem estudante perspectivas motivadoras na sua \textbf{trajetória} escolar, ao procurar adequar e \textbf{atender} diferentes demandas a partir da significação e valorização de históricos locais.
Nesse sentido, o \textbf{Currículo} Paulista etapa Ensino Médio demanda novas formas de organização dos tempos e espaços escolares, da gestão do \textbf{currículo}, das metodologias e da formação dos professores.
Nesse processo, a flexibilização, materializada nos \textbf{itinerários} \textbf{formativos}, apresenta formas de propiciar novas experiências de organização e estrutura do ensino nas escolas, que deverão ser incorporadas nas suas Propostas Pedagógicas e/ou Regimentos Escolares.

% matched lemmas: adolescente, alinhar, atender, capítulo, concorrer, currículo, curso, destinar, diretriz, escolaridade, formativo, frequentar, habilitação, homologar, itinerário, legislação, masculino, oferta, política, portaria, profissional, qualidade, qualificação, rede, sintonizar, trabalhador, trajetória, técnico
\end{textsample}
